% Preámbulo común de los informes del TP5 (señal 1 y señal 2), con el estilo de los trabajos anteriores.
\documentclass[12pt]{article}
\usepackage[utf8]{inputenc}
\usepackage[T1]{fontenc}
\usepackage{graphicx}
\usepackage[spanish,es-tabla]{babel}
\usepackage{amsmath}
\usepackage{amssymb}
\usepackage[a4paper,margin=1in]{geometry}
\usepackage{subcaption}
\usepackage{mathtools}
\usepackage[table,xcdraw]{xcolor}
\usepackage[style=ieee, backend=bibtex]{biblatex}
\usepackage{enumitem}
\usepackage{array}
\usepackage{longtable}
\usepackage{multirow}
\usepackage{booktabs}
\usepackage{csquotes}
\usepackage{tikz}
\usetikzlibrary{arrows.meta, positioning, shapes.geometric, babel}
\usepackage[hidelinks]{hyperref}
\usepackage{float}
\usepackage{setspace}
\onehalfspacing
\graphicspath{{Figuras/}}

\addbibresource{bibliografia.bib}
\hypersetup{colorlinks=true, linkcolor=blue, citecolor=blue, urlcolor=blue, filecolor=blue}

% Portada: \portada{integrantes (tabla)}{señal}
\newcommand{\portada}[2]{%
\begin{titlepage}
\thispagestyle{empty}
\begin{tikzpicture}[remember picture,overlay]
    \node[opacity=0.2,anchor=center] at (current page.center)
        {\includegraphics[width=0.9\paperwidth]{decoracion_a.png}};
\end{tikzpicture}
\begin{center}
    \includegraphics[width=15cm]{unc_fcefyn_logo.png} \\ [0.5cm]
    \LARGE Universidad Nacional de Córdoba \\ [0.5cm]
    \large Facultad de Ciencias Exactas, Físicas y Naturales \\
    \vspace{40pt}
    \large Sistemas de Radiocomunicación \\ [0.2cm]
    \large Trabajo Práctico N.° 5 \\
    \vspace{50pt}
    \LARGE Radar \\[0.3cm]
    \large Procesamiento de la señal #2 \\
    \vspace{50pt}
    \begin{table}[!h]
        \centering
        \begin{tabular}{ll}
            \multicolumn{1}{c}{Nombre} & \multicolumn{1}{c}{DNI} \\
            #1
        \end{tabular}
    \end{table}
    \vspace{20pt}
    \begin{table}[!h]
        \centering
        \begin{tabular}{ll}
            \multicolumn{1}{c}{Docentes} & Esp. Ing. Liendo, Carlos Guillermo \\
                                         & Ing. Dadam, Federico Tomás \\
        \end{tabular}
    \end{table}
    \vfill
        {\small Córdoba, República Argentina\\ Año 2026}
\end{center}
\end{titlepage}
}
